\documentclass[11pt]{article}
% Standalone preamble for web-only supplementary notes.
\usepackage[T1]{fontenc}
\usepackage{amsmath,amssymb,amsthm}
\usepackage{hyperref}
\usepackage{cleveref}
\DeclareMathOperator{\Map}{Map}
\DeclareMathOperator{\Fun}{Fun}
\newcommand{\iso}{\xrightarrow{\sim}}
\newtheorem{lemma}{Supplementary Lemma}
\newtheorem{corollary}[lemma]{Supplementary Corollary}
\crefname{lemma}{Supplementary Lemma}{Supplementary Lemmas}
\Crefname{lemma}{Supplementary Lemma}{Supplementary Lemmas}
\crefname{corollary}{Supplementary Corollary}{Supplementary Corollaries}
\Crefname{corollary}{Supplementary Corollary}{Supplementary Corollaries}

\begin{document}
\section*{Endpoint coordinates for slice categories}

The proof describing \href{slice-categories.html}{functors into a slice}
changes from a defining endpoint pullback to a cylinder description.
Here are the coordinate calculations used in that change.
Assume the Chapter 1 structures, including functor categories and pullbacks,
and the walking morphism. All constructions in this note are absolute.

Write
$\chi_{X,Y,C}\colon\Fun(X,\Fun(Y,C))\iso\Fun(X\times Y,C)$
for the exponential comparison, and $u^*$ for precomposition by $u$.
For $\psi\colon Y\to C$, use the direct absolute slice: the pullback of
endpoint evaluation
\[
 \Fun([1]\times Y,C)\longrightarrow\Fun(Y,C)\times\Fun(Y,C)
\]
along $C\to\Fun(Y,C)^2$, $c\mapsto(\operatorname{const}_c,\psi)$.
This is the absolute slice implementation linked below. Its comparison
with the separately chosen relative slice over $\star$ is not exported
by these declarations.

\begin{lemma}
\label[lemma]{lem:Supplementary_Slice_Exponential_Restriction}
For $u\colon Y\to Z$ there is an identification of functors
\[
 \chi_{X,Y,C}\,\Fun(X,u^*)
 = (\operatorname{id}_X\times u)^*\,\chi_{X,Z,C}.
\]
\end{lemma}
\begin{proof}[Proof idea]
Uncurry both sides. Each becomes double evaluation after a change of
product coordinates. Compare the two reassociation routes by their three
projections: the first two coordinates are unchanged and the last is
postcomposed with $u$. The exponential reflection rule gives the claimed
identification. In particular this applies to either endpoint insertion
$Y\to[1]\times Y$.
\end{proof}

\begin{lemma}
\label[lemma]{lem:Supplementary_Slice_Constants}
Under $\chi_{X,Y,C}$, postcomposition by the constant-family functor
$C\to\Fun(Y,C)$ becomes
$\operatorname{pr}_X^*\colon\Fun(X,C)\to\Fun(X\times Y,C)$.
Postcomposition by the constant functor $C\to\Fun(Y,C)$ with value
$\psi$ becomes the constant functor with value $\psi\operatorname{pr}_Y$.
These are identifications of functors on $\Fun(X,C)$.
\end{lemma}
\begin{proof}[Proof idea]
Double evaluation of the first construction forgets the $Y$ coordinate
and evaluates the variable functor at $X$. Double evaluation of the second
forgets the variable functor and the $X$ coordinate and evaluates $\psi$
at $Y$. The product-projection and constant-family computation laws
identify these maps with the uncurrying of the two asserted functors.
Reflect through the exponential equivalence.
\end{proof}

\begin{lemma}
\label[lemma]{lem:Supplementary_Slice_Cylinder_Endpoints}
The equivalence
\[
 r\colon X\times([1]\times Y)\iso(X\times Y)\times[1]
\]
given by symmetry and reassociation preserves endpoint insertion:
\[
 r\bigl(\operatorname{id}_X\times(e,\operatorname{id}_Y)\bigr)
 = (\operatorname{id}_{X\times Y},e).
\]
There is also a comparison for the chosen inverse of $r$ with the
insertion on the left.
\end{lemma}
\begin{proof}[Proof idea]
Compute the $X,Y,[1]$ projections. The first two are the projections of
$X\times Y$; the third is the constant endpoint $e$.
Pair their comparisons to obtain the displayed identification.
For the inverse comparison, compose with $r$, use its retraction
comparison, and reflect through the equivalence $r$.
\end{proof}

\begin{corollary}
\label[corollary]{cor:Supplementary_Tested_Slice}
Put $A=\Fun(X,C)$, $Z=\Fun(X\times Y,C)$ and
$W=\Fun((X\times Y)\times[1],C)$.
There is a pullback square with corners
\[
 \begin{matrix}
 \Fun(X,C_{/\psi})&\longrightarrow&W\\
 \big\downarrow&&\big\downarrow\\
 A&\longrightarrow&Z\times Z,
 \end{matrix}
\]
where the right map is endpoint restriction, the lower map sends $f$ to
$(f\operatorname{pr}_X,\psi\operatorname{pr}_Y)$, and the left map is
postcomposition with the slice projection. Its commutativity
identification is the one transported from the defining slice square
by the preceding comparisons.
\end{corollary}
\begin{proof}[Proof idea]
Apply $\Fun(X,-)$ to the defining slice pullback. It remains a pullback.
Use the exponential and product comparisons on its corners.
\Cref{lem:Supplementary_Slice_Exponential_Restriction,lem:Supplementary_Slice_Cylinder_Endpoints}
identify the upper cospan map with endpoint restriction;
\Cref{lem:Supplementary_Slice_Constants} identifies the lower one with
the displayed pair. These computations give a comparison of cospans,
including its specified squares. Transport the cone through that comparison.
Pullback invariance supplies the assertion with the transported matching.

The lower map factors through
$\Fun(X,C)\times\Fun(Y,C)$ by $f\mapsto(f,\psi)$.
With the mapping-out join axiom, nesting the two pullbacks and using
the mapping-out description of the join gives the square in the book. The coordinate calculation above is
what makes that nesting apply to the specified slice diagram.
\end{proof}
\end{document}
